\documentclass[10pt]{article}
\usepackage[margin=20mm]{geometry}
\usepackage{amsmath,amssymb,mathtools,fontspec,unicode-math,xurl,hyperref}
\IfFontExistsTF{DejaVu Serif}{\setmainfont{DejaVu Serif}}{\setmainfont{Latin Modern Roman}}
\setmathfont{latinmodern-math.otf}
\setlength{\parindent}{0pt}\setlength{\parskip}{6pt}\sloppy
\title{Technical verification of KPZ fluctuations in the planar stochastic heat equation}
\author{AI-assisted technical review}
\date{}
\begin{document}
\maketitle
Abstract: The general scaling statement in Theorem 6 loses a factor of the time-scale parameter in the noise strength. An exact occupation-coordinate calculation and a bounded Gaussian test show that its asserted coupled mean limit is false when that parameter differs from one. The natural corrected strength is beta times rho. This is a Major issue in the stated general theorem; the counterexample does not refute the rho equals one specialization or its KPZ interpretation. A separate Major proof defect occurs in the high-probability truncation used for the central convergence argument: the printed uncentered bridge cutoff is empty at large drift. The fixed-cutoff repair and its limits are stated explicitly. Other auxiliary formulas also require correction: classical divergence normalization and domain, a polymer heat-kernel denominator, Gaussian pair counting, and local-time scaling and truncation bounds. Explicit counterexamples and the limits of the repairs are given below. Scientific adjudication remains pending.\par
\section*{Exact source and scope}
The source is Jeremy Quastel, Alejandro Ramírez and Bálint Virág, KPZ fluctuations in the planar stochastic heat equation, M25-0164, arXiv:2210.13607v3, 10 April 2024, 69 pages [source]. Physical PDF and printed page numbers agree on inspected pages. The assignment associates the article with DOI 10.1215/00127094-2024-0058 (https://doi.org/10.1215/00127094-2024-0058); no identity with a journal-final version is asserted.\par
The review covers Sections 1–7: the \(L^1\) divergence definition, finite projections and uniqueness, randomized shifts and their convergence, occupation moments, intersection local times, truncated exponential estimates, and assembly of the main results. Native scientific helpers examined disjoint sections; the lead checked the findings and the central normalization calculations. Coverage, exact file hashes, and actual page-view records accompany this report. Calculations are analytic. No unavailable author program is treated as a mathematical defect. The bibliography was not independently re-proved; the one external normalization comparison used below is identified precisely.\par
\section*{The general scaling theorem has the wrong noise strength}
Problem 1 (Major: Theorem 6, pages 8–9 and 66).  With the parameters and common-noise coupling stated in Theorem 6, the limiting noise strength is \(\beta\rho\), whereas the theorem prints \(\beta\).\par
Equation reference: eq:rho. Let the target time interval have length \(s>0\). The planar bridge runs for \(N\rho s\), and the scaled basis is \[e_{N,j}(x,y)=N^{-1}e_j(N^{-1/2}x,N^{-3/2}y).\] At time \(N\rho r\), its coordinates have the representation \[B_N(N\rho r)=
 \bigl(N^{1/2}b_{\rm sp}(t+r),\,
       N^{3/2}(t+r)+N^{1/2}b_{\rm ti}(t+r)\bigr),
 \quad 0\le r\le s,\] where the spatial bridge has variance parameter \(\nu\rho\), and the centered transverse bridge has variance parameter \(\rho\). Consequently its occupation coordinate, including the noise coefficient, is exactly \[\begin{aligned}
 m_{N,j}
 &=\beta\int_0^{N\rho s} e_{N,j}(B_N(u))\,du \notag\\
 &=\beta\rho\int_0^s
 e_j\bigl(b_{\rm sp}(t+r),t+r+N^{-1}b_{\rm ti}(t+r)\bigr)\,dr.
 \end{aligned}\] The change of variable on page 9 omits \(\rho\). For bounded continuous \(e_j\), the limiting coordinate is thus \(\beta\rho\int_0^s e_j(b_{\rm sp}(t+r),t+r)\,dr\). The deterministic heat-kernel normalization in the theorem does not remove this factor.\par
\subsection*{A bounded-test contradiction to the printed \(L^1\) limit}
Take \(\rho=2\), \(\nu=\beta=s=1\), and both spatial endpoints and the initial time equal to zero. Choose as the first orthonormal basis function \[e_1(x,r)=\sqrt{2/\pi}\exp(-x^2-r^2).\] It has \(L^2(\mathbb R^2)\) norm one and can be extended to an orthonormal basis of bounded continuous functions, by Gram–Schmidt applied to a dense family containing it. Write \(M=\|e_1\|_\infty=\sqrt{2/\pi}\) and \[m=\int_0^1 e_1(b_{\rm sp}(r),r)\,dr.\] Every continuous bridge path gives \(0<m\le M\), and [eq:rho] (\#eq:rho) gives \(m_{N,1}\to2m\), with \(|m_{N,1}|\le2M\).\par
Divide the rescaled solution by its deterministic limiting spatial heat kernel, and denote the resulting mean-one partition function by \(Z_N\). The randomized-shift identity, (65) or (85) in the source, gives for \(a=\pi/(8M)\) \[\mathbb E_P[Z_N\sin(a\xi_1)]
 =e^{-a^2/2}\mathbb E_Q[\sin(a m_{N,1})]
 \longrightarrow e^{-a^2/2}\mathbb E_Q[\sin(2am)].\] The partition function \(Z_*\) of the printed target \(z_{2,1}\) instead satisfies \[\mathbb E_P[Z_*\sin(a\xi_1)]
 =e^{-a^2/2}\mathbb E_Q[\sin(am)].\] Since \(0<am\le\pi/8\), the first limiting expectation is strictly larger. A bounded test must have convergent expectations under \(L^1(P)\) convergence. This proves a contradiction to the exact printed theorem, without assuming uniform second moments or invoking a merely formal chaos expansion.\par
\subsection*{Repair, boundaries, and effect}
Replace the target \(z_{\nu\rho,\beta}\) by \(z_{\nu\rho,\beta\rho}\), with the same correction in the polymer coupling, or replace the planar noise coefficient by \(\beta/\rho\). This is the necessary algebraic repair. The deterministic prefactor \(N\sqrt{2\pi\rho s}\exp(N^2s/(2\rho))\) is consistent with the remaining spatial heat kernel. The counterexample uses strictly positive admissible parameters, not a degenerate endpoint. At \(\rho=1\) the missing factor disappears, so Theorem 5 is not refuted. A full convergence proof still requires the local-time and truncated-exponential inputs discussed in Section 5 (\#sec:ilt); correcting one coefficient alone is not presented as a new proof of every statement of Theorem 6.\par
\section*{Classical divergence and mild-equation normalization}
Problem 2 (Minor: Section 2.5, pages 19–23).  The tensor conventions in (36)–(37) do not give the claimed comparison with Definition 8. The displayed divergence norm and full-domain description on page 21 are also false.\par
The paper uses unnormalized symmetrization \(\operatorname{Sym}g=\sum_{\sigma}g\circ\sigma\). In one dimension (36) therefore gives \(I_k(e^{\otimes k})=H_k(\xi)/k!\). For \(G=H_2(\xi)=\xi^2-1\), its kernel in that convention is \(g=2e^{\otimes2}\). Formula (37) yields \[I_3\bigl(\operatorname{Sym}(2e^{\otimes3})\bigr)=2H_3(\xi),\] whereas Definition 8 and Gaussian integration by parts give \(\delta G=\xi G-G'=H_3(\xi)\). This discrepancy survives the admissible bounded test \(F(\xi)=\sin\xi\): \[\mathbb E[H_3(\xi)\sin\xi]
 =\mathbb E[H_2(\xi)\cos\xi]=-e^{-1/2},\] while the printed candidate gives twice this value. All terms are finite polynomials with every finite moment.\par
There is a separate obstruction to the asserted norm equality. On a two-coordinate Gaussian space let \(G=(\xi_2,-\xi_1)\). Then \(\delta G=0\), while the claimed weighted chaos norm is \((1+1)\mathbb E\|G\|^2=4\). Symmetrization can have a kernel. The weighted condition (41) is sufficient, but is not the full domain. Indeed, for \(k\ge1\), put \[v_k=
 \frac{(H_{k-1}(\xi_1)\xi_2,\,-H_k(\xi_1))}
 {\sqrt{(k-1)!+k!}} .\] These mutually orthogonal vector fields have norm one, lie in chaos \(k\), and satisfy \(\delta v_k=0\). Thus \(G=\sum_{k\ge1}k^{-1}v_k\) is in \(L^2\) and in the classical divergence domain with divergence zero, but \[\sum_{k\ge1}(k+1)\mathbb E\|\pi_kG\|^2
 =\sum_{k\ge1}\frac{k+1}{k^2}=\infty.\] Closed Gaussian adjoint duality, or the convergent divergence chaos series, justifies the domain assertion in this example.\par
A consistent repair uses averaged symmetrization and standard Wiener integrals throughout. For kernels \(f_k\in H^{\odot k}\otimes H\), the actual condition is \(\sum_{k\ge0}(k+1)!\|\widetilde f_k\|^2<\infty\), and \(\mathbb E|\delta G|^2\) equals that sum. If the paper’s rescaled \(I_k\) convention is retained, (37) needs a factor \(1/k!\), and the occupation-coordinate formulas must be adjusted consistently.\par
\subsection*{Function integrability versus coordinate integrability}
Lemma 25 must also distinguish Definition 8 from Definition 10. The deterministic function \(g(x)=(1+|x|)^{-3/4}\) on \(\mathbb R\) has \(\int g^2=4\), so its classical Wiener integral exists, but \(\int |g|=\infty\). It does not satisfy Definition 10’s finite absolute-mass hypothesis. The coordinate interpretation can be correct; the function interpretation needs that additional hypothesis.\par
The needed PDE comparison can be obtained without the false tensor norm identity. Finite Hermite integration by parts gives the Gaussian adjoint relation. Correct \(L^2\) chaos summation extends it to the divergence domain. For predictable step integrands, integration by parts in an independent future increment identifies this divergence with the Itô integral; step approximation and Itô isometry extend the identity on the predictable \(L^2\) class.\par
\subsection*{The missing heat factor and the missing endpoint integral}
In Example 28, already at first chaos, the mild kernel is \[g_1(y;t,x)=\int_0^t p(x-y,t-r)p(y,r)\,dr.\] A probability bridge’s expected occupation density is \(g_1(y;t,x)/p(x,t)\). Thus (46) needs the deterministic factor \(p(x,t)\), consistently with the chronological heat products in (47).\par
Equation (49), pages 23–24, incorrectly equates a fixed-endpoint heat convolution with the full \(L^1\) norm. With \(t=1\), \(v(y,r)=p(y,r)\) and \(x=0\), its sides are \(p(0,1)=1/(2\pi)\) and \(1\). The sufficient estimate is instead, with \[A(x,t)=\mathbb E\int_{\mathbb R^2}\int_0^t
 p(x-y,t-r)|u(y,r)|\,dr\,dy,\] \[\int_0^T\int_{\mathbb R^2} A(x,t)\,dx\,dt
 =\int_0^T(T-r)\mathbb E\|u(\cdot,r)\|_1\,dr
 \le T\|u\|_{L^1(dx\,dr\,dP)}.\] Tonelli therefore gives the required finite absolute mass and coordinate integrability at almost every endpoint. This repairs the Fubini steps used in Propositions 29–30 and 33 under their stated \(L^1\) hypothesis.\par
\section*{Randomized shifts and polymer masses}
Problem 3 (Minor: Proposition 39(c), Example 47, and Lemma 87).  Three auxiliary normalizations need correction. Their effects are distinct from the missing time-scale factor in Theorem 6.\par
\subsection*{Gaussian pair counting}
The exponent in Proposition 39(c), page 30, sums over ordered pairs \(j\ne\ell\). It needs a factor \(1/2\), or the sum should be over \(j<\ell\). For one deterministic coordinate \(m_1=1\) and \(k=2\), \(Z=e^{\xi-1/2}\) has \(\mathbb EZ^2=e\), while the printed formula gives \(e^2\). Gaussian completion of the square gives exactly \[\mathbb EZ_n^k=
 \mathbb E_{Q^{\otimes k}}
 \exp\left\{\sum_{j<\ell}\sum_{i=1}^n m_i^{(j)}m_i^{(\ell)}\right\}.\] The separately printed second-moment identity, Proposition 39(b), is correct and is the identity needed in the later \(L^2\) criteria.\par
\subsection*{The circle illustration}
Example 47, page 35, combines covariance \(2\log|e^{i\omega}-e^{i\omega'}|^{-1}\) with the exponential \(\exp(\gamma\zeta-\gamma^2\mathbb E\zeta^2/2)\), but states the subcritical range \(\gamma<2\). With that covariance and exponential, the range is \(\gamma<1\). In the explicitly checked comparison [Remy], page 2, (1.2)–(1.3), the range \(<2\) uses the coefficient \(\gamma/2\) in front of the same field. Substituting \(2\gamma\) gives the correction. The source’s \(L^2\) threshold \(1/\sqrt2\) is consistent with its actual exponential. This is an ancillary example, not an input to the Brownian-bridge construction.\par
\subsection*{A deterministic denominator in Lemma 87}
The reference bridge law \(Q\) is a probability measure, and the random measure \(M\) defined by (85) has \(\mathbb E_P M=Q\), hence \(\mathbb E_P M(\Omega)=1\). In Lemma 87, page 65, the displayed finite-time density should be \[\mathbb E_M\varphi
 =\frac{1}{q(x_{k+1}-x_0,s_{k+1}-s_0)}
 \int_{\mathbb R^k}\varphi(x_1,\ldots,x_k)
 \prod_{i=0}^k z(x_i,s_i;x_{i+1},s_{i+1})\,dx_1\cdots dx_k .\] The source omits the denominator. Its own (118) is correct: multiplying the segment partition functions under the probability bridge produces the bridge density \(\prod_i q_i/q_{\rm total}\).\par
For an exact contradiction take \(k=1\), intermediate time \(1/2\), endpoints \((0,0),(0,1)\), and \(\varphi=1\). The expectation of the left side is \(1\). Independence of the two time-noise slabs, \(\mathbb Ez=q\), and the heat semigroup identity give expectation \(q(0,1)=1/\sqrt{2\pi}\) on the printed right side. Restoring the denominator also gives the familiar normalized polymer density \(\prod_i z_i/z_{\rm total}\). It does not invalidate the existence of the random measure.\par
\subsection*{The meaning of the convergence criterion}
Proposition 52(2) concerns convergence of nonnegative martingale limits. Its limiting cutoff condition should not be silently read as existence of a mean-one randomized-shift density for every approximating model. This distinction matters when invoking the random measures in Proposition 60.\par
For clarity, a concrete boundary example is available. Let \(U\) be uniform on \([0,1]\), independently of a fair infinite binary path. Index the nonempty tree nodes \(v\), and put \(h_v=2\,1_{\{v\text{ is a prefix}\}}\). Two independent paths share a prefix of length \(L\) with \(\mathbb P(L\ge n)=2^{-n}\); their intersection sum is \(4L\), with all finite polynomial moments. The depth-\(n\) partition martingale \(W_n\) satisfies \[\mathbb EW_n^{1/2}\le(\sqrt2e^{-1/2})^n\longrightarrow0,\] by subadditivity of the square root, so its limit has lost all mass. Let \(\varepsilon_\ell=(\ell+1)^{-1}\) and \(m_\ell=1_{\{U\le\varepsilon_\ell\}}h\), with limit \(m=0\). The coordinates converge almost surely, the second intersection moments tend to zero, and for each \(p<1\) the cutoffs \(A_\ell=\{U>\varepsilon_\ell\}\) eventually have probability \(>p\) and zero intersection on their product. Nevertheless, every martingale limit is the constant \(1-\varepsilon_\ell\), not a mean-one density. These constants do converge to \(1\) in \(L^1\); thus this example does not refute Proposition 52’s convergence assertion. Existence of the individual GMCs must be supplied where they are used. The Brownian applications have separate uniform-integrability arguments.\par
\section*{Intersection-local-time estimates}
Problem 4 (Major: the large-drift truncation used in Theorem 6).  The high-probability truncation needed in the main convergence proof is not established by the printed argument. Several local-time formulas also require correction; weaker corrected bounds suffice for the stated limiting application.\par
\subsection*{Occupation mass is quadratic in the time change}
In Lemma 76, page 53, \(\widehat X=tX\circ L^{-1}\). Bilinearity gives \[\alpha(\widehat X,\widehat X')
 =\frac{t^2}{|\det L|}\alpha(X,X'),\] not \(t/|\det L|\). For the deterministic path \(B(r)=r\) on \([0,1]\), in dimension one, take \(L=1\) and \(t=2\). The original occupation density is \(1_{[0,1]}\), the new one is \(2\,1_{[0,1]}\), and their self-intersection values are \(1\) and \(4\), not \(1\) and \(2\). This exact finite-density example satisfies the definition without any singular limiting issue.\par
In Proposition 74(2), page 52, duplication of the gamma function reduces the exact displayed moment to \[\mathbb E\alpha^k=t^{k/2}\Gamma(1+k/2).\] Stirling’s formula gives order \(\sqrt{k}\,(tk/(2e))^{k/2}\), not the printed upper bound with \(1/\sqrt{k}\). The all-exponential-moments conclusion nevertheless survives: \(\mathbb E\alpha^{2m}=t^m m!\) gives \(\mathbb E\cosh(\gamma\alpha)\le e^{\gamma^2t}\).\par
For Lemma 78, page 55, the bridge bound \(e^{10\gamma}\) should be \(e^{10\gamma t}\), and the drift second-moment numerator should be \(ct^2\), not \(ct\). These are genuine failures for unrestricted \(t\). At zero drift, planar intersection local time scales as \(\alpha_t\stackrel d=t\alpha_1\). Thus \(\mathbb E\alpha_t^2=t^2\mathbb E\alpha_1^2\) cannot be bounded by \(ct\) for all \(t\). For bridges, take \(\gamma=1/t\); the permitted condition \(\gamma t=1\) holds, but the left side is the fixed number \(\mathbb Ee^{\alpha_1}>1\), whereas \(e^{10/t}\to1\). The proof at \(t=1\) and correct Brownian scaling yield the repaired forms. Lemma 82 already prints the correct \(t^2\) numerator.\par
\subsection*{Proposition 77: the quantity being limited}
The definition (102) omits a square: the proof and (103) concern \(\phi(\nu,y)=\mathbb E[\alpha(B^{\nu,y},(B^{\nu,y})')^2]\). At \(\nu=0,y=0,t=1\), Proposition 74 gives \(\mathbb E\alpha=\sqrt{\pi}/2\) and \(\mathbb E\alpha^2=1\), so the distinction is substantive. The Gaussian integral on page 54 also has \(2\pi\nu^{-2}\sqrt{\det C}\) in the denominator where \(2\pi\nu^2\sqrt{\det C}\) is required. After these corrections, the remaining step is a singular approximate identity limit near the paired time diagonals; it needs control of endpoint and collision neighborhoods, not only pointwise passage under a Fourier integral.\par
Here is a completion for the intended second moment, independent of Propositions 84–85. Put \(G(s)_{ij}=\min(s_i,s_j)-s_i s_j\), \(C=G(s)+G(t)\), \(\Delta=s-t\), and \(Q=\Delta^\top C^{-1}\Delta\). The nonnegative four-time kernel is \[K_{\nu,y}(s,t)
 =\frac{\exp[-Q/(2\nu^2)-y^2Q/2]}
 {(2\pi)^2\nu^2\det C},\qquad
 M_\nu(y)=\int_{[0,1]^4}K_{\nu,y}.\] In particular \(K_{\nu,y}\le K_{\nu,0}\). Set \(N=1/\nu\). The verified \(k=2\) density bound from Propositions 68–69 gives, for motions of length \(\ell\) with shared drift \((0,N)\), \[\mathbb E\alpha^2\le\frac{C\ell^2}{1+N^2\ell}.\] Stopping each first coordinate at their midpoint gives the stronger bound with factor \(\exp[-|x-x'|^2/(8\ell)]\): before either hit the paths cannot intersect; after the hit apply the strong Markov property, extend the remaining segment, and use the uniform starting-point bound. This uses only the second-moment part of Lemma 82. No exponential stopping estimate is needed.\par
On a bridge prefix ending at \(v<1\), the pair Radon–Nikodym bound relative to drifted motions is \((1-v)^{-2}\). Splitting at \(1/2\), using time reversal and moment-density Cauchy–Schwarz for mixed halves, proves \(\sup_{\nu\le1}M_\nu(0)<\infty\). For endpoint windows \(E_a=[0,a]\cup[1-a,1]\), the same bound gives \(\mathbb E\alpha(X_{E_a},X'_{E_a})^2\le Ca\) after the anisotropic factor \(N^2\). Applying Cauchy–Schwarz twice shows that all four-time terms having at least one endpoint time in \(E_a\) contribute at most \(Ca^{1/4}\).\par
Since \(C_{ii}\le1/2\), \(Q\ge2\Delta_i^2\). Off \(\max_i|s_i-t_i|\le\eta\), split the exponential to obtain \[K_{\nu,y}\le2e^{-\eta^2/(2\nu^2)}K_{\sqrt2\nu,0}.\] The integral there tends to zero. On the remaining interior region, if \(|s_1-s_2|\le h\) and \(\max_i|s_i-t_i|\le h/4\), all four times lie in one of \(O(h^{-1})\) windows of length \(O(h)\). Under the prefix comparison, condition on both starting positions at time \(u\ge a/2\). Their horizontal difference has law \(D_u\sim N(0,2u)\), so \[\mathbb Ee^{-D_u^2/(8\ell)}
 =(1+u/(2\ell))^{-1/2}\le\sqrt{2\ell/u}.\] The anisotropically scaled second moment of each window is at most \(C_a\ell^{3/2}\); summing gives \(C_a\sqrt h\).\par
After excluding these regions, \(C\) is uniformly positive definite. The substitution \(t=s+\nu w\) has an integrable Gaussian majorant, so dominated convergence gives density \(1/(2\pi\sqrt{\det(2G(s))})\). First send \(\nu\to0\), then \(h\to0\), and then \(a\to0\). The full limiting integral is \[\frac1{2\pi}\int_{0<u<v<1}
 [u(v-u)(1-v)]^{-1/2}\,du\,dv=1.\] Thus the intended \(M_\nu(y)\to1\) is verified for every fixed \(y\), after the two printed corrections. The companion proof record gives the cutoff details; neither the literal first-moment definition nor the later exponential bound is validated by this calculation.\par
\subsection*{Proposition 84 cannot have the printed uniform prefactor}
In Proposition 84, page 61, fix \(N=1\) and both starting points at zero. Let \(A_r\) be its stated centered-strip event and let \(m=\mathbb E\alpha(B,B')>0\). Gaussian maximal tails give \(\mathbb P((A_r\times A_r)^c)\le C e^{-r^2/8}\), and integrability gives \(\mathbb E[\alpha1_{A_r\times A_r}]\to m\). Set \(\gamma=r^{-2}\), which satisfies \(\gamma\le\pi/(8r)\) for all sufficiently large \(r\). Since \(e^x\ge1+x\), \[\mathbb E[e^{\gamma\alpha}1_{A_r\times A_r}]
 \ge1-Ce^{-r^2/8}+\frac{m}{2r^2},
 \qquad
 e^{c\gamma^2}=1+O(r^{-4}).\] No fixed \(c\) makes the printed inequality true for all \(r\). The factor \(2\) in the paper’s (115) cannot be absorbed into \(e^{c\gamma^2}\) uniformly as \(\gamma\downarrow0\).\par
A sufficient repair retains a prefactor \(C_0>1\). The interval covering also needs slack: \(|s-t|\le r/N\) on the strip does not ensure that a pair belongs to one of two half-shifted grids of windows of length \(r/N\). Windows of length \(2r/N\), with shift \(r/N\), supply such a cover, at the cost of a smaller universal constant in the admissible bound \(\gamma r<c_0\). The Gaussian weight must also use the actual block duration: at duration \(r/N\) it is \(\exp[-NY_i^2/(8r)]\), rather than the \(\exp[-Y_i^2/(8r)]\) printed near (115). Only the former becomes a fixed-parameter Gaussian weight after normalizing the increments to a standard random walk. These changes do not justify retaining every printed numerical constant.\par
\subsection*{What the large-drift application actually needs}
The event \(H\) in Proposition 85 uses the uncentered norm from (107). For a bridge from \(0\) to \((0,N)\), \(\|B\|_\kappa\ge N\), so that event is empty when \(N>r/2\). It cannot be the high-probability cutoff needed for the limit.\par
On page 66, the proof of Theorem 6 explicitly invokes Proposition 85 to supply the hypothesis of Proposition 52(2): for every \(p<1\), there must be path sets \(A_N\) with \(Q(A_N)>p\) and a finite limiting upper bound on the truncated exponential moments. That page fixes \(\nu=\rho=\beta=s=1\) and \(x=t=0\). The printed empty bridge event cannot provide those sets. After centering, the covering and exponential estimates still require the corrections and conditional argument below. Thus the displayed proof does not establish a necessary uniform control in the central convergence argument, including its standard-parameter specialization. This is a Major proof defect, not a counterexample to that limit. This operative dependency determines the assessment of this group; the other local coefficient and scaling errors are not each thereby classified as Major.\par
Use instead \[H_r=\{\|B-(0,N)\,\cdot\|_\kappa\le r/2\}.\] Its probability can be made arbitrarily close to one independently of \(N\). This centered event is also what implies the increment conditions \(G_i\) used on page 62.\par
After choosing fixed \(r,\gamma,\kappa\), select a finite \(m\) large enough that the corrected short-block bound applies, and then \(N\) large enough for the separation of the blocks. The conditional block bound needs an additional argument because each increment cutoff is centered at its own starting position. For two such independent path measures \(X,Y\), with respective cutoffs \(G,G'\), moment-density Cauchy–Schwarz at each order and then Cauchy–Schwarz in the nonnegative exponential series give \[\mathbb E[e^{\theta\alpha(X,Y)}1_G1_{G'}]
 \le\left(
 \mathbb E[e^{\theta\alpha(X,X')}1_G1_{\widetilde G}]\,
 \mathbb E[e^{\theta\alpha(Y,Y')}1_{G'}1_{\widetilde G'}]
 \right)^{1/2}.\] The zeroth term is the product of cutoff probabilities. For every positive order, the moment densities are in \(L^2\) by the untruncated Brownian moment bounds. The self terms are identical up to a common translation, because the increment cutoff is translation covariant. Each can therefore use the corrected strip bound. No independent translation of the two paths inside a cross intersection local time is being assumed. The retained prefactor costs \(C_0^m\), which is finite and independent of \(N\). The resulting bound is \[\sup_{N\ge N_0(r,\gamma,\kappa)}
 \mathbb E[e^{\gamma N\alpha(B,B')}1_{H_r\times H_r}]
 <\infty .\] For motions, one may take \(m\) larger than a fixed multiple of \(1+(\gamma r)^{1/\kappa}\) and \(N\) larger than a fixed multiple of \(rm\). For bridges, the overlapping two-thirds prefixes and their pair Radon–Nikodym factor \(9\), followed by time reversal and Cauchy–Schwarz, preserve a finite bound depending on the fixed cutoff parameters. This weaker bound is the one needed by the cutoff convergence criterion. The displayed choice \(m=\lceil(16\gamma r/\pi)^{1/\kappa}\rceil\) does not by itself establish the printed cubic threshold in \(N\) for every \(\kappa<1/2\). That numerical threshold is not certified here; no counterexample to it is claimed.\par
Two further proof completions are bounded. In Lemma 83, the chosen \(a=4\gamma e^{\lambda/2}/\sqrt{k}\) need not be at most one under the stated \(k\ge\gamma^2\). For \(k\ge16e^\lambda\gamma^2\) the argument applies; below that threshold the trivial estimate \(\exp(\gamma\sqrt{k})\le\exp(4e^{\lambda/2}\gamma^2)\) is covered by the claimed right side. In Proposition 79, the grid-measurable occupation cutoff must remain in the conditional estimate through (108)–(111). Dropping it before invoking its bound on the number of pairs is not justified. Keeping it gives the finite exponential bound used by Corollary 81.\par
\section*{Verified dependence and limits of the conclusion}
The identification in Lemma 67 need not rely on pointwise convergence of raw Dirichlet Fourier sums from continuity alone. Heat-kernel regularizations of the convolution have a second-moment inner product controlled by the \(L^2\) occupation moment densities. Strong continuity of heat convolution makes them Cauchy in \(L^2\); the defining rectangle-test identity identifies the limit with \(\alpha\). For a realized finite convolution measure whose density is continuous near zero, the positive heat approximate identity converges to that density value. These two limits prove the identification used in the paper.\par
The finite-dimensional construction in Section 3 can be checked directly from Definition 8. On a finite time interval, \(|m(r)|\le r(\sum_{j\le n}\|e_j\|_\infty^2)^{1/2}\); the Wick exponential and its needed Gaussian derivatives therefore have uniform integrable bounds. Bridge disintegration and the correct Tonelli estimate prove existence. Fourier tests in the Gaussian coordinates reduce uniqueness to an \(L^1\) heat equation with bounded potential. An \(L^1\) Gronwall argument proves uniqueness; using rational Fourier parameters first and then continuity handles the common-null-set issue.\par
For Proposition 33, \(\mathbb E|u_n(x,t)|\le P_t|\varsigma|(x)\). Pointwise \(L^1(P)\) convergence is consequently dominated in the full space–time product \(L^1\) norm. Conditional expectations of the limiting kernel agree with the finite projected kernels, so Corollary 14 applies. These arguments do not depend on the false classical norm equality.\par
The randomized-shift identity follows from finite Gaussian completion of the square. Mean-one martingale uniform integrability, the cutoff Cauchy argument of Proposition 45, the coordinate moment bounds of Propositions 48–50, and the \(L^2\) part of Proposition 52 were checked. For an already existing random measure \(M\), finite-coordinate approximations satisfy \(M_n(F)=\mathbb E_P[M(F)\mid\mathcal F_n]\) for every bounded path test \(F\). Conditional-expectation convergence and tightness from \(\mathbb EM_n=Q\) give weak convergence in probability on the Polish path space. This supplies the approximation step without assuming that a convergence criterion creates every individual GMC.\par
Two Major assessments are supported: the explicit false general parameter limit in Theorem 6, and the missing high-probability uniform exponential control in its printed tightness proof. The latter also affects the standard-parameter derivation on page 66; the verified centered fixed-cutoff repair restores the needed hypothesis, without proving the stronger printed constants or cubic threshold. It is not a counterexample to the standard-parameter limit. The remaining auxiliary issues have the specified local effects. The exact sharp constants and all formulations of the repaired auxiliary statements are not claimed as independently optimized. No journal-final comparison, author-code reproduction, or comprehensive later-literature search was performed. The local evidence and the necessary corrections are tied to the exact supplied version.\par
9 J. Quastel, A. Ramírez and B. Virág, KPZ fluctuations in the planar stochastic heat equation, arXiv:2210.13607v3, 10 April 2024, supplied 69-page manuscript. G. Remy, The Fyodorov–Bouchaud formula and Liouville conformal field theory, arXiv:1710.06897v3, 4 November 2019, p. 2, equations (1.2)–(1.3) and the following parameter range. https://arxiv.org/pdf/1710.06897.\par
\end{document}
